% ---------------------------------------------------------------------------
% AI-Sociology arXiv package — shared preamble
% Engine: XeLaTeX (required: Cyrillic in article_ru.tex, Unicode glyphs)
% This file carries no article prose.
% ---------------------------------------------------------------------------

\usepackage{fontspec}

% --- fonts -----------------------------------------------------------------
% Font cascade: first available wins. Every candidate covers Latin and
% Cyrillic, so the package compiles on a minimal TeX Live as well as on the
% arXiv build host.
\newcommand{\pickmain}[1]{\IfFontExistsTF{#1}{\setmainfont{#1}\global\let\pickdone\relax}{}}
\newcommand{\pickmono}[1]{\IfFontExistsTF{#1}{\setmonofont{#1}[Scale=MatchLowercase]\global\let\monodone\relax}{}}

% Candidates are ordered by typographic preference AMONG fonts that carry both
% Latin and Cyrillic. TeX Gyre Termes is last because some TeX Live builds ship
% it without Cyrillic coverage (verified locally: 17 missing glyphs), which
% would silently drop the entire Russian edition from the PDF.
\ifdefined\pickdone\else\pickmain{Libertinus Serif}\fi
\ifdefined\pickdone\else\pickmain{Liberation Serif}\fi
\ifdefined\pickdone\else\pickmain{DejaVu Serif}\fi
\ifdefined\pickdone\else\pickmain{FreeSerif}\fi
\ifdefined\pickdone\else\pickmain{TeX Gyre Termes}\fi

\ifdefined\monodone\else\pickmono{DejaVu Sans Mono}\fi
\ifdefined\monodone\else\pickmono{Liberation Mono}\fi
\ifdefined\monodone\else\pickmono{Latin Modern Mono}\fi

% --- languages -------------------------------------------------------------
% English is the main language; the Russian and Dutch editions are wrapped in
% their own language environments so that hyphenation and spacing follow the
% edition rather than the host document.
% NOTE: a minimal TeX Live without texlive-lang-cyrillic / -european loads no
% Russian or Dutch hyphenation patterns and warns accordingly. The arXiv build
% host carries the full set, so those editions hyphenate correctly there; the
% only local symptom is additional overfull boxes inside those two editions.
\usepackage{polyglossia}
\setmainlanguage{english}
\setotherlanguage{russian}
\setotherlanguage{dutch}

% Second guard: bind the Russian edition to a family verified to carry
% Cyrillic, so the Russian text survives even if the main font does not.
\newcommand{\pickru}[1]{\IfFontExistsTF{#1}{\newfontfamily\russianfont{#1}\global\let\rudone\relax}{}}
\ifdefined\rudone\else\pickru{Libertinus Serif}\fi
\ifdefined\rudone\else\pickru{Liberation Serif}\fi
\ifdefined\rudone\else\pickru{DejaVu Serif}\fi
\ifdefined\rudone\else\pickru{FreeSerif}\fi

\usepackage{amssymb}
\usepackage{tabularx}
\usepackage{xltabular}
\usepackage{booktabs}
\usepackage{longtable}
\usepackage{microtype}
\usepackage[hidelinks,unicode]{hyperref}

% --- glyph macros ----------------------------------------------------------
% The frozen masters use these characters. They are routed through macros so
% the printed character never depends on one particular font being installed
% (verified locally: DejaVu Serif carries neither U+2713 nor U+26A0).
% Semantics are unchanged; only the glyph shape is normalised.
\newcommand{\gcheck}{\checkmark}                      % ✓ verified reference
\newcommand{\gwarn}{\ensuremath{\triangle}}           % ⚠ provisional reference
\newcommand{\gharr}{\ensuremath{\leftrightarrow}}     % ↔
\newcommand{\grarr}{\ensuremath{\rightarrow}}         % →
\newcommand{\gtimes}{\ensuremath{\times}}             % ×
\newcommand{\gsect}{\S}                               % §

% --- section cross-references ---------------------------------------------
% \secref{4.5.4} prints "§4.5.4" exactly as the frozen text does, and links to
% that section when the label exists in the current edition. All 294 uses in
% the English master resolve; none is dangling.
\newcommand{\EditionPrefix}{en}
\makeatletter
\newcommand{\secref}[1]{%
  \edef\sr@target{sec:\EditionPrefix:#1}%
  \ifcsname r@\sr@target\endcsname
    \expandafter\hyperref\expandafter[\sr@target]{\S#1}%
  \else
    \S#1%
  \fi}
\makeatother

% --- structures used by the converted editions -----------------------------
\newcounter{dummysec}

% Front-matter contributor block: the frozen masters carry this as a
% preformatted block. Content is reproduced word for word; only the
% typesetting is changed.
% \obeylines turns each source newline into a line break, so the block needs
% no explicit \\ and the preformatted layout of the master is preserved.
\newenvironment{frontmatterblock}
  {\par\medskip\begingroup\small\raggedright
   \setlength{\parindent}{0pt}\setlength{\parskip}{0pt}
   \obeylines\obeyspaces}
  {\endgroup\par\medskip}

% Starred sectioning is used throughout, because the section numbers are part
% of the frozen text (and the English master has no §8.4). Automatic numbering
% would silently renumber and break the printed numbers that the article's own
% 294 cross-references depend on.
\setcounter{secnumdepth}{0}
\setcounter{tocdepth}{2}

% Tables in these editions are prose-heavy; allow a little more tolerance than
% the default so that narrow cells do not produce spurious overfull boxes.
% A small global emergency stretch absorbs the few lines whose long hyphenated
% compounds and identifiers cannot otherwise be broken. Typographic only.
\emergencystretch=3em
\setlength{\tabcolsep}{4pt}
\renewcommand{\arraystretch}{1.12}

% --- bibliography ----------------------------------------------------------
\usepackage[numbers,sort&compress]{natbib}
\bibliographystyle{plainnat}
